\documentclass[11pt,a4paper]{article}

\usepackage[margin=2.2cm]{geometry}
\usepackage[T1]{fontenc}
\usepackage[utf8]{inputenc}
\usepackage{lmodern}
\usepackage{amsmath,amssymb}
\usepackage{enumitem}
\usepackage{parskip}

\begin{document}

\section*{Suggested Answers --- trial\_exam\_6}

\subsection*{Section 1 --- Multiple Choice (Q1--Q20)}
\begin{itemize}
    \item Q1 B
    \item Q2 A
    \item Q3 A
    \item Q4 A
    \item Q5 A
    \item Q6 A
    \item Q7 B
    \item Q8 B
    \item Q9 C
    \item Q10 A
    \item Q11 A
    \item Q12 C
    \item Q13 C
    \item Q14 A
    \item Q15 B
    \item Q16 B
    \item Q17 B
    \item Q18 B
    \item Q19 A
    \item Q20 A
\end{itemize}

\subsection*{Section 2 --- Short Questions (Q1--Q6)}
\begin{enumerate}[label=\arabic*.]
    \item If symbol duration is $T_s$, then $R_s = 1/T_s$.
    \item The decision samples occur at the receiver's symbol timing instants (i.e., aligned to symbol boundaries determined by clock/timing recovery).
    \item $\eta$ is the spectral efficiency; with $\eta=\frac{R_b}{B}$, higher $\eta$ means more information bits per second are carried per Hz.
    \item \textbf{Duplexing:} how uplink and downlink share resources (time or frequency division).
    \item \textbf{Resource allocation goal/challenge:} allocate spectrum/time/power to meet performance targets; a key challenge is managing interference (and meeting QoS/scalability constraints).
    \item Shorter symbols are more sensitive because higher symbol rate generally implies larger required bandwidth (more noise collected) and the channel can cause more overlap between symbols (worse ISI/distortion).
\end{enumerate}

\subsection*{Section 3 --- Long / Applied Questions}

\subsubsection*{Question 1 --- Link Budget and Feasibility}
Given: $f=700$ MHz, $P_t=10$ dBm, $G_t=0$ dBi, $G_r=2$ dBi, $d=1.8$ km, $N=-102$ dBm, $\gamma_{\min}=7$ dB.

\begin{enumerate}[label=\alph*.]
    \item \textbf{Free-space path loss:}
    \[
    L_{\mathrm{FSPL}}=32.44+20\log_{10}(700)+20\log_{10}(1.8)\approx 94.45\ \mathrm{dB}.
    \]
    \item \textbf{Received power:}
    \[
    P_r=P_t+G_t+G_r-L_{\mathrm{FSPL}}=10+0+2-94.45\approx -82.45\ \mathrm{dBm}.
    \]
    \item \textbf{Minimum required received power:}
    \[
    P_{\min}=N+\gamma_{\min}=-102+7=-95\ \mathrm{dBm}.
    \]
    \item \textbf{Feasible?} yes, since $P_r\approx -82.45 > -95$ dBm.
    \item (Optional margin) $\approx 12.55$ dB.
    
\end{enumerate}

\subsubsection*{Question 2 --- Indoor Path Loss}
Given: distance loss $75$ dB, two concrete walls ($2\times 8$ dB), door $5$ dB, glass $3$ dB, $P_t=16$ dBm, $G_t=G_r=0$ dBi, sensitivity $-82$ dBm.

\begin{enumerate}[label=\alph*.]
    \item Total path loss:
    \[
    L_{\text{total}}=75+2\cdot 8+5+3=99\ \mathrm{dB}.
    \]
    \item Received power:
    \[
    P_r=P_t-L_{\text{total}}=16-99=-83\ \mathrm{dBm}.
    \]
    \item Link works? $-83$ dBm is $1$ dB below $-82$ dBm, so \textbf{no}.
    \item Two improvements (examples): increase transmit power; use higher-gain antennas/better placement; reduce obstacles (or add relay).
\end{enumerate}

\subsubsection*{Question 3 --- Frequency Reuse}
Given reuse factor $N=2$.

\begin{enumerate}[label=\alph*.]
    \item Fraction of spectrum per cell:
    \[
    \frac{1}{2}=0.5=50\%.
    \]
    \item For $N=6$:
    \[
    \frac{1}{6}\approx 0.167\ (16.7\%).
    \]
    \item Impact of $N$:
    \begin{itemize}
        \item Smaller $N$ $\Rightarrow$ more spectrum per cell but more co-channel interference.
        \item Larger $N$ $\Rightarrow$ less co-channel interference but less spectrum per cell (lower nominal capacity).
    \end{itemize}
    
\end{enumerate}

\subsubsection*{Question 4 --- Multi-hop Routing}
Feasible links: $1\to 2,\ 1\to 3,\ 2\to 3,\ 2\to 4,\ 3\to 4,\ 1\to 4$.

Energies: $E_{12}=2.6,\ E_{13}=3.1,\ E_{23}=1.2,\ E_{24}=2.0,\ E_{34}=2.4,\ E_{14}=5.9$.

\begin{enumerate}[label=\alph*.]
    \item All feasible routes from node $1$ to node $4$:
    \begin{itemize}
        \item Route A: $1\to 4$
        \item Route B: $1\to 2\to 4$
        \item Route C: $1\to 3\to 4$
        \item Route D: $1\to 2\to 3\to 4$
    \end{itemize}
    \item Total energy per route:
    \[
    E_A=E_{14}=5.9
    \]
    \[
    E_B=E_{12}+E_{24}=2.6+2.0=4.6
    \]
    \[
    E_C=E_{13}+E_{34}=3.1+2.4=5.5
    \]
    \[
    E_D=E_{12}+E_{23}+E_{34}=2.6+1.2+2.4=6.2
    \]
    \item Minimum-energy route: Route B with $E_{\min}=4.6$.
    \item Advantage: improved feasibility/coverage with smaller per-hop distances. Disadvantage: more delay and overhead due to more hops.
\end{enumerate}

\subsubsection*{Question 5 --- Modulation and Throughput}
Bandwidth: $B=110$ kHz $=1.10\times 10^5$ Hz.

\begin{enumerate}[label=\alph*.]
    \item Shannon capacity at SNR $=4$ dB:
    \[
    C=B\log_2(1+\mathrm{SNR})\approx 1.993\times 10^5\ \text{bit/s}\ (\approx 199.35\ \text{kbps}).
    \]
    \item Highest usable scheme at $4$ dB: scheme-3 (threshold $\ge 4$ dB). Throughput $\approx 1.5$ bit/s/Hz.
    \[
    R\approx 1.5\cdot B=1.5\cdot 110000=1.65\times 10^5\ \text{bit/s}\ (165\ \text{kbps}).
    \]
    \item SNR drop from scheme-4 to scheme-3:
    scheme-4 is usable at SNR $\ge 8$ dB. At exactly $8$ dB there is effectively no margin before switching, so
    \[
    \Delta \approx 0\ \text{dB}.
    \]
\end{enumerate}

\subsubsection*{Question 6 --- Coverage Planning}
\begin{enumerate}[label=\alph*.]
    \item \textbf{Single-site:} place the base station at the highest/central location to improve line-of-sight or diffraction paths over the hill.
    \item \textbf{Two-site (base + relay):} use a relay positioned on the far-side of the hill (or on a ridge) so the shadowed town receives a stronger hop.
    \item \textbf{Compare:} two-site typically improves coverage and reliability but increases deployment cost and complexity.
    \item \textbf{Most important propagation effects:} hill blockage/shadowing, diffraction, and free-space path loss (plus reflections/multipath near buildings).
\end{enumerate}

\end{document}

